\subsection{实验结果图表分析}

本节汇总模型对比、混淆矩阵、训练过程、固定案例消融和压力测试结果。
图表数据来自训练输出、模型评估文件和实验日志；
与前文表格中的数值保持一致。
训练数据按照 60\%、15\% 和 25\% 划分训练集、验证集和独立测试集；
独立测试集规模为 2280 条。

不同情感分类模型的独立测试集结果见图~\ref{fig:model-performance}。

\begin{figure}[H]
  \centering
  \includegraphics[width=0.92\textwidth]{model_performance.png}
  \caption{独立测试集情感分类模型对比}
  \label{fig:model-performance}
\end{figure}

图~\ref{fig:model-performance} 将传统基线和最终 BERT 模型置于同一坐标系中。
评价指标为 \mbox{Accuracy} 和 \mbox{Macro-F1}。
multilingual BERT 的 Accuracy 为 0.7461，Macro-F1 为 0.7457；
分别比最佳传统模型 Logistic Regression 提高约 5.57 和 5.51 个百分点。
这说明上下文表示对课程评价中的转折、否定和混合态度有实际帮助。

BERT 在三类情感上的具体错误分布见图~\ref{fig:bert-confusion-matrix}。

\begin{figure}[H]
  \centering
  \includegraphics[width=0.61\textwidth]{bert_confusion_matrix.png}
  \caption{multilingual BERT 独立测试集混淆矩阵}
  \label{fig:bert-confusion-matrix}
\end{figure}

图~\ref{fig:bert-confusion-matrix} 显示，积极类别识别效果最好。
中性类别有 228 条被误判为消极、72 条被误判为积极；
这是最主要的错误来源。
课程评价中的建议型表达、轻度抱怨和正负混合表达仍容易与单一极性混淆。
因此，不能只依赖总体指标判断模型效果。

BERT 微调过程见图~\ref{fig:bert-training-process}。

\begin{figure}[H]
  \centering
  \begin{subfigure}[t]{0.48\textwidth}
    \centering
    \includegraphics[width=\textwidth]{bert_training_loss.png}
    \caption{训练损失}
    \label{fig:bert-training-loss}
  \end{subfigure}
  \hfill
  \begin{subfigure}[t]{0.48\textwidth}
    \centering
    \includegraphics[width=\textwidth]{bert_validation_metrics.png}
    \caption{验证集指标}
    \label{fig:bert-validation-metrics}
  \end{subfigure}
  \caption{multilingual BERT 微调过程}
  \label{fig:bert-training-process}
\end{figure}

图~\ref{fig:bert-training-process}(a) 显示训练损失在两个 Epoch 内总体下降。
图~\ref{fig:bert-training-process}(b) 显示验证集 Accuracy 从 0.7003 提升至 0.7230；
Macro-F1 从 0.7039 提升至 0.7225。
第二轮训练仍带来稳定增益；
但验证损失未同步明显下降，因此当前训练轮数保持为 2，以控制过拟合风险和训练成本。

固定案例消融结果见图~\ref{fig:ablation-visualization}。

\begin{figure}[H]
  \centering
  \includegraphics[width=0.90\textwidth]{ablation_results.png}
  \caption{固定案例消融结果}
  \label{fig:ablation-visualization}
\end{figure}

固定案例主要用于课堂演示流程验证，不作为泛化能力指标。
项目分别运行纯规则、TF-IDF、BERT 和完整混合流程。
图~\ref{fig:ablation-visualization} 显示，TF-IDF 单模块通过 7/12；
BERT 单模块通过 11/12；
完整流程通过 12/12。
纯规则同样在这组固定案例上达到满分，原因是案例集中包含较多规则能够直接覆盖的典型表达。
因此，项目进一步使用未参与规则开发的独立压力测试检验鲁棒性。

独立压力测试结果见图~\ref{fig:stress-test-visualization}。

\begin{figure}[H]
  \centering
  \begin{subfigure}[t]{0.72\textwidth}
    \centering
    \includegraphics[width=\textwidth]{stress_test_overall.png}
    \caption{整体 Accuracy 与 Macro-F1}
    \label{fig:stress-test-overall}
  \end{subfigure}

  \vspace{0.5em}
  \begin{subfigure}[t]{\textwidth}
    \centering
    \includegraphics[width=\textwidth]{stress_test_categories.png}
    \caption{分类别通过数}
    \label{fig:stress-test-categories}
  \end{subfigure}
  \caption{独立压力测试结果}
  \label{fig:stress-test-visualization}
\end{figure}

压力测试共 48 条，覆盖缩写、拼写、否定、讽刺、隐含抱怨、中英混合、长文本和混合情感；
每类 6 条。
压力测试不是随机测试集，而是面向困难表达构造的挑战集。
其结果主要反映复杂边界场景下的鲁棒性，不应与 2280 条常规独立测试集指标直接比较。
图~\ref{fig:stress-test-visualization} 表明，混合流程在压力测试上的 Accuracy 和 Macro-F1
均高于三个单独模块。
其优势主要集中在缩写、拼写错误、否定、中英混合和长文本场景。
另一方面，讽刺仅通过 1/6，隐含抱怨仅通过 3/6；
这些类型仍是当前系统的主要短板。
